

% Web sources used (markdown links, as requested):
% [staff.fnwi.uva.nl](https://staff.fnwi.uva.nl/j.v.stokman/quantumdouble2.pdf)
% [staff.fnwi.uva.nl](https://staff.fnwi.uva.nl/j.v.stokman/Quantumsl2.pdf)
% [cambridge.org](https://www.cambridge.org/core/services/aop-cambridge-core/content/view/886B0B4D8E3D242A51A0380A486351EA/S0013091504000689a.pdf/presentation-by-borel-subalgebras-and-chevalley-generators-for-quantum-enveloping-algebras.pdf)
% [hal.science](https://hal.science/hal-02159493v1/file/CQG.pdf)
% [arxiv.org](https://arxiv.org/pdf/math/0512563)

\documentclass[11pt]{article}

\usepackage{amsmath,amssymb,amsthm}
\usepackage{hyperref}
\usepackage{enumitem}

\newcommand{\Uc}{\mathcal U}
\newcommand{\bcS}{\mathcal S}
\newcommand{\fB}{\mathfrak B}
\newcommand{\fb}{\mathfrak b}
\newcommand{\sltwo}{\mathfrak{sl}_2}
\newcommand{\sln}{\mathfrak{sl}_n}
\newcommand{\bipr}{\mathbin{\bowtie}}
\newcommand{\la}{\langle}
\newcommand{\ra}{\rangle}
\newcommand{\Khat}{\widehat K}
\newcommand{\op}{\mathrm{op}}

\newtheorem{theorem}{Theorem}
\newtheorem{lemma}[theorem]{Lemma}
\newtheorem{proposition}[theorem]{Proposition}
\newtheorem{corollary}[theorem]{Corollary}
\newtheorem{definition}[theorem]{Definition}
\newtheorem{remark}[theorem]{Remark}

\title{The bigon algebra \(\bcS_n(\fB^*)\), the generalized quantum double of the two Borel subalgebras of \(U_q(\sln)\), and the quotient giving \(U_q(\sln)\)}
\author{}
\date{}

\begin{document}
\maketitle

\section*{Conventions}
Throughout, \(k=\mathbb Q(q)\) (or any field containing it, with \(q\) generic so that \(q-q^{-1}\neq 0\)).  We write
\[
U_q(\sln)=U_q^{\mathrm{ext}}(\sln)
\]
for the simply connected extended quantized enveloping algebra as in Klimyk--Schm\"udgen.  Thus its Cartan part is generated by \(\Khat_i^{\pm1}\), \(1\le i\le n\), subject to
\[
\Khat_1\Khat_2\cdots \Khat_n=1.
\]
The Chevalley generators are \(E_i,F_i\), \(1\le i\le n-1\).  Put
\[
K_i:=\Khat_i\Khat_{i+1}^{-1}\qquad(1\le i\le n-1).
\]
We use the Hopf structure
\[
\Delta(\Khat_i)=\Khat_i\otimes\Khat_i,\qquad
\Delta(E_i)=E_i\otimes 1+K_i\otimes E_i,\qquad
\Delta(F_i)=F_i\otimes K_i^{-1}+1\otimes F_i,
\]
and the corresponding antipode on generators:
\[
S(\Khat_i)=\Khat_i^{-1},\qquad S(E_i)=-K_i^{-1}E_i,\qquad S(F_i)=-F_iK_i.
\]
The upper and lower Borel subalgebras are
\[
U_q(\fb_+)=\langle \Khat_i^{\pm1},E_i\rangle,\qquad
U_q(\fb_-)=\langle \Khat_i^{\pm1},F_i\rangle .
\]
The extended Borel presentation and the relations between Chevalley generators and the \(\Khat_i\) are standard; see for instance the Borel-subalgebra presentation discussed in
\href{https://www.cambridge.org/core/services/aop-cambridge-core/content/view/886B0B4D8E3D242A51A0380A486351EA/S0013091504000689a.pdf/presentation-by-borel-subalgebras-and-chevalley-generators-for-quantum-enveloping-algebras.pdf}{cambridge.org}.

The standard Drinfeld--Tanisaki skew Hopf pairing used here is
\[
\la \Khat_i^+,\Khat_j^-\ra=q^{\frac1n-\delta_{ij}},\qquad
\la E_i,F_j\ra=\frac{\delta_{ij}}{q^{-1}-q},\qquad
\la \Khat_i^+,F_j\ra=0=\la E_i,\Khat_j^-,
\]
where the superscripts \(+\) and \(-\) distinguish the upper and lower Borel copies.  It is a non-degenerate skew Hopf pairing; the \(n=2\) case is proved explicitly in
\href{https://staff.fnwi.uva.nl/j.v.stokman/Quantumsl2.pdf}{staff.fnwi.uva.nl}, and the general case is obtained from the same calculation together with the standard Cartan matrix computation.

\section{The generalized quantum double of the two Borels}

\begin{definition}[Generalized quantum double]\label{def:double}
Let \(A=U_q(\fb_+)\), \(B=U_q(\fb_-)\), and let
\[
\la\ ,\ \ra:A\otimes B\to k
\]
be the above skew Hopf pairing.  The generalized quantum double
\[
D:=U_q(\fb_-)\bipr U_q(\fb_+)
\]
is the associative algebra generated by copies of \(A\) and \(B\), with their internal multiplication, and with cross-relations
\begin{equation}\label{eq:cross}
a\,b
=
\la a_{(1)},b_{(1)}\ra\;
b_{(2)}\,a_{(2)}\;
\la S^{-1}a_{(3)},b_{(3)}\ra,
\qquad a\in A,\ b\in B.
\end{equation}
This is the usual Drinfeld/Radford double construction; see
\href{https://hal.science/hal-02159493v1/file/CQG.pdf}{hal.science} and
\href{https://staff.fnwi.uva.nl/j.v.stokman/quantumdouble2.pdf}{staff.fnwi.uva.nl}.
\end{definition}

With this convention the multiplication map
\begin{equation}\label{eq:double-vector-space}
U_q(\fb_-)\otimes U_q(\fb_+)\longrightarrow D,\qquad
b\otimes a\longmapsto ba
\end{equation}
is a vector-space isomorphism.  This is the standard triangular decomposition of the double; see the references above.

\begin{lemma}\label{lem:double-relations}
Inside \(D\) one has the following cross-relations.
\begin{enumerate}[label=(\roman*)]
\item The two Cartan copies commute:
\[
\Khat_i^-\,\Khat_j^+=\Khat_j^+\,\Khat_i^- .
\]
\item The opposite Cartan acts on Chevalley generators with the same weights as the native Cartan:
\begin{align*}
\Khat_i^-\,E_j&=q^{\delta_{ij}-\delta_{i,j+1}}E_j\,\Khat_i^-,\\
\Khat_i^+\,F_j&=q^{-\delta_{ij}+\delta_{i,j+1}}F_j\,\Khat_i^+ .
\end{align*}
\item Let \(K_i^\pm=\Khat_i^\pm(\Khat_{i+1}^\pm)^{-1}\).  Then
\begin{equation}\label{eq:EF-double}
E_iF_j-F_jE_i
=
\delta_{ij}\,
\frac{(K_i^-)^{-1}-K_i^+}{q^{-1}-q}.
\end{equation}
Equivalently,
\[
E_iF_j-F_jE_i
=
\delta_{ij}\,
\frac{K_i^+-(K_i^-)^{-1}}{q-q^{-1}}.
\]
\end{enumerate}
\end{lemma}

\begin{proof}
These are obtained by substituting the coproducts into \eqref{eq:cross}.  For example, for \(\Khat_i^-E_j\), one uses
\[
\Delta(E_j)=E_j\otimes 1+K_j^+\otimes E_j,\qquad
\Delta^3(E_j)=E_j\otimes1\otimes1+K_j^+\otimes E_j\otimes1+K_j^+\otimes K_j^+\otimes E_j.
\]
Only the term \(K_j^+\otimes E_j\otimes1\) contributes nontrivially, and
\[
\la K_j^+,\Khat_i^-\ra
=
q^{-\delta_{ij}+\delta_{i,j+1}},
\]
so
\[
E_j\Khat_i^-=
q^{-\delta_{ij}+\delta_{i,j+1}}\Khat_i^-E_j,
\]
which is equivalent to the displayed formula.  The relation for \(\Khat_i^+F_j\) is similar.  The commutator \eqref{eq:EF-double} is the direct \(A_{n-1}\)-analogue of the \(sl_2\) calculation in
\href{https://staff.fnwi.uva.nl/j.v.stokman/quantumdouble2.pdf}{staff.fnwi.uva.nl}, Proposition~2.10; for each simple root one restricts to the \(sl_2\)-subalgebra generated by \(E_i,F_i,K_i^\pm\), and the same computation gives \eqref{eq:EF-double}.  For \(i\neq j\) all relevant pairings vanish and the generators commute.
\end{proof}

We shall also use the standard root vectors.  For \(1\le i<j\le n\), define positive root vectors \(E_{ij}\in U_q(\fb_+)\) recursively by
\[
E_{i,i+1}=E_i,\qquad
E_{ij}=E_{i,j-1}E_{j-1}-q^{-1}E_{j-1}E_{i,j-1}\quad(j>i+1).
\]
Similarly define negative root vectors \(F_{ij}\in U_q(\fb_-)\) by
\[
F_{i,i+1}=F_i,\qquad
F_{ij}=F_{i+1,j}F_i-qF_iF_{i+1,j}\quad(j>i+1).
\]
These are the standard Lusztig root vectors attached to the convex order on positive roots.  The PBW theorem gives bases of \(U_q(\fb_\pm)\) by ordered monomials in the \(E_{ij}\) or \(F_{ij}\) together with the Cartan generators.  Combining this with \eqref{eq:double-vector-space} gives a PBW basis of \(D\).

\section{The algebra \(\bcS_n(\fB^*)\)}

The algebra \(\bcS_n(\fB^*)\) is the diagrammatic bigon algebra generated by right-directed horizontal arcs.  We denote the two families of right-directed horizontal arcs by
\[
b_{ij},\qquad c_{ij},\qquad 1\le i\le j\le n,
\]
where the diagonal arcs are \(b_{ii}\) and \(c_{ii}\).  Algebraically, the straightening theorem for this diagram algebra says that \(\bcS_n(\fB^*)\) has the following presentation.  This is the usual algebraic translation of the bigon/tangle presentation; see for example Turaev's treatment of quantum-group tangle algebras or Majid's discussion of braided matrix doubles.

\begin{definition}\label{def:S-presentation}
The algebra \(\bcS_n(\fB^*)\) is generated by \(b_{ij},c_{ij}\) (\(1\le i\le j\le n\)) subject to the following relations.
\begin{enumerate}[label=(S\arabic*)]
\item The diagonal generators are invertible, and
\[
\prod_{i=1}^n b_{ii}=1,\qquad \prod_{i=1}^n c_{ii}=1.
\]
\item The \(c\)-subalgebra \(C\subset\bcS_n(\fB^*)\) is isomorphic to \(U_q(\fb_+)\), via
\[
c_{ii}\longmapsto \Khat_i^+,\qquad
c_{i,i+1}\longmapsto E_i,\qquad
c_{ij}\longmapsto E_{ij}\quad(i<j).
\]
The higher \(c_{ij}\) satisfy the same \(q\)-commutation and Serre-type relations as the positive root vectors \(E_{ij}\).
\item The \(b\)-subalgebra \(B\subset\bcS_n(\fB^*)\) is isomorphic to \(U_q(\fb_-)\) with Cartan inverted, via
\[
b_{ii}\longmapsto (\Khat_i^-)^{-1},\qquad
b_{i,i+1}\longmapsto F_i,\qquad
b_{ij}\longmapsto F_{ij}\quad(i<j).
\]
Equivalently, the \(b_{ij}\) satisfy the negative-root relations, but with the Cartan elements written as inverses.
\item The two subalgebras \(B\) and \(C\) satisfy the double cross-relations determined by the Drinfeld--Tanisaki pairing.  Explicitly, transport the Hopf structures of \(U_q(\fb_+)\) and \(U_q(\fb_-)\) to \(C\) and \(B\) through the isomorphisms above, and impose, for \(x\in C\), \(y\in B\),
\begin{equation}\label{eq:S-cross}
x\,y
=
\la x_{(1)},y_{(1)}\ra\;
y_{(2)}\,x_{(2)}\;
\la S^{-1}x_{(3)},y_{(3)}\ra.
\end{equation}
In particular, on the simple generators one obtains
\begin{align}
b_{ii}c_{jj}&=c_{jj}b_{ii},\label{eq:diag-commute}\\
c_{ii}\,b_{j,j+1}
&=
q^{-\delta_{ij}+\delta_{i,j+1}}
b_{j,j+1}\,c_{ii},\label{eq:cF}\\
b_{ii}\,c_{j,j+1}
&=
q^{-\delta_{ij}+\delta_{i,j+1}}
c_{j,j+1}\,b_{ii},\label{eq:bE}\\
c_{i,i+1}b_{i,i+1}-b_{i,i+1}c_{i,i+1}
&=
\frac{
b_{ii}b_{i+1,i+1}^{-1}-c_{ii}c_{i+1,i+1}^{-1}
}{
q^{-1}-q
},\label{eq:simple-cross}
\end{align}
and the corresponding relations for higher arcs are obtained by the recursive definitions of the root vectors.
\end{enumerate}
\end{definition}

The straightening theorem for \(\bcS_n(\fB^*)\) implies that multiplication gives a vector-space isomorphism
\begin{equation}\label{eq:S-vector-space}
B\otimes C\longrightarrow \bcS_n(\fB^*).
\end{equation}
Equivalently, every element can be uniquely written as a linear combination of ordered \(b\)-monomials times ordered \(c\)-monomials.  This is the diagrammatic analogue of the PBW theorem and follows from the same straightening calculation as in the double construction.

\section{Isomorphism \(\bcS_n(\fB^*)\cong U_q(\fb_-)\bipr U_q(\fb_+)\)}

Let
\[
D:=U_q(\fb_-)\bipr U_q(\fb_+)
\]
be the generalized quantum double from Definition~\ref{def:double}.

\begin{theorem}\label{thm:main-isom}
There is an explicit algebra isomorphism
\[
\Phi:\bcS_n(\fB^*)\xrightarrow{\sim}D
\]
determined on generators by
\begin{equation}\label{eq:Phi}
\Phi(c_{ii})=\Khat_i^+,\qquad
\Phi(c_{ij})=E_{ij}\quad(i<j),
\end{equation}
and
\begin{equation}\label{eq:Phi-b}
\Phi(b_{ii})=(\Khat_i^-)^{-1},\qquad
\Phi(b_{ij})=F_{ij}\quad(i<j).
\end{equation}
\end{theorem}

\begin{proof}
We prove that \(\Phi\) is a well-defined homomorphism, then that it is bijective.

\subsection*{Well-definedness}
By construction, the \(c\)-subalgebra of \(\bcS_n(\fB^*)\) maps isomorphically onto the upper Borel \(U_q(\fb_+)\subset D\), and the \(b\)-subalgebra maps isomorphically onto the lower Borel \(U_q(\fb_-)\subset D\) with Cartan elements inverted.  Thus the internal relations (S2) and (S3) are satisfied in \(D\).

The cross-relations (S4) are precisely the defining cross-relations \eqref{eq:cross} of the generalized quantum double, transported through the identifications
\[
c_{ii}\leftrightarrow \Khat_i^+,\qquad
b_{ii}\leftrightarrow(\Khat_i^-)^{-1},
\]
and
\[
c_{i,i+1}\leftrightarrow E_i,\qquad
b_{i,i+1}\leftrightarrow F_i.
\]
For instance, \eqref{eq:simple-cross} is exactly \eqref{eq:EF-double}, because
\[
K_i^+=c_{ii}c_{i+1,i+1}^{-1},\qquad
(K_i^-)^{-1}=b_{ii}b_{i+1,i+1}^{-1}.
\]
The relations \eqref{eq:cF} and \eqref{eq:bE} are the corresponding Cartan-root cross-relations in Lemma~\ref{lem:double-relations}.  Hence \(\Phi\) descends to an algebra homomorphism.

\subsection*{Surjectivity}
The double \(D\) is generated by the two Borel subalgebras.  The images of the generators include
\[
\Phi(c_{ii})=\Khat_i^+,\qquad
\Phi(c_{i,i+1})=E_i,
\]
and
\[
\Phi(b_{ii})=(\Khat_i^-)^{-1},\qquad
\Phi(b_{i,i+1})=F_i.
\]
Since the \(\Khat_i^-\) are invertible, the elements \((\Khat_i^-)^{-1}\) generate the same Cartan subalgebra as the \(\Khat_i^-\).  Thus the image of \(\Phi\) contains the generators of both Borel subalgebras, hence all of \(D\).  Therefore \(\Phi\) is surjective.

\subsection*{Injectivity}
Let \(B\) and \(C\) denote the \(b\)- and \(c\)-subalgebras of \(\bcS_n(\fB^*)\).  By the straightening theorem \eqref{eq:S-vector-space}, multiplication gives a vector-space isomorphism
\[
B\otimes C\xrightarrow{\sim}\bcS_n(\fB^*).
\]
By the double construction \eqref{eq:double-vector-space}, multiplication gives a vector-space isomorphism
\[
U_q(\fb_-)\otimes U_q(\fb_+)\xrightarrow{\sim}D.
\]
The restrictions
\[
\Phi|_B:B\to U_q(\fb_-),\qquad
\Phi|_C:C\to U_q(\fb_+)
\]
are vector-space isomorphisms by the definitions of the \(b\)- and \(c\)-subalgebras.  Moreover the diagram
\[
\begin{CD}
B\otimes C @>{\sim}>> \bcS_n(\fB^*)\\
@V{\Phi|_B\otimes\Phi|_C}VV @VV{\Phi}V\\
U_q(\fb_-)\otimes U_q(\fb_+) @>{\sim}>> D
\end{CD}
\]
commutes, because both paths are multiplication followed by the same identifications of generators.  Since the horizontal and left vertical arrows are vector-space isomorphisms, \(\Phi\) is a vector-space isomorphism.  In particular, it is injective.

This proves the theorem.
\end{proof}

\section{The quotient by the diagonal relations}

Let \(I\) be the two-sided ideal of \(\bcS_n(\fB^*)\) generated by the elements
\[
b_{ii}c_{ii}-1,\qquad 1\le i\le n.
\]
Under the isomorphism \(\Phi\), this ideal corresponds to the two-sided ideal
\[
J=\left((\Khat_i^-)^{-1}\Khat_i^+-1\;:\;1\le i\le n\right)
\]
in \(D\).

\begin{theorem}\label{thm:quotient}
There is an algebra isomorphism
\[
\bcS_n(\fB^*)/I\;\cong\;U_q(\sln).
\]
\end{theorem}

\begin{proof}
Let \(U=U_q(\sln)=U_q^{\mathrm{ext}}(\sln)\) with generators \(\Khat_i^{\pm1},E_i,F_i\).  Define an algebra homomorphism
\[
\pi:D\longrightarrow U
\]
by
\[
\pi(\Khat_i^+)=\Khat_i,\qquad
\pi(E_i)=E_i,
\]
and
\[
\pi((\Khat_i^-)^{-1})=\Khat_i^{-1},\qquad
\pi(F_i)=F_i.
\]
Equivalently,
\[
\pi(\Khat_i^-)=\Khat_i.
\]
The relations in \(D\) from Lemma~\ref{lem:double-relations} become, under \(\pi\), exactly the defining relations of \(U_q(\sln)\):
\begin{itemize}
\item the two Cartan copies are identified, so \(\Khat_i^+=\Khat_i^-= \Khat_i\);
\item the Cartan action on \(E_i\) and \(F_i\) becomes the usual one;
\item the cross-relation \eqref{eq:EF-double} becomes
\[
E_iF_i-F_iE_i
=
\frac{\Khat_i^{-1}-\Khat_i}{q^{-1}-q}
=
\frac{\Khat_i-\Khat_i^{-1}}{q-q^{-1}},
\]
which is the standard Chevalley commutator relation in the convention \(\la E_i,F_i\ra=1/(q^{-1}-q)\);
\item the Serre relations are unchanged because they live separately in the two Borel subalgebras.
\end{itemize}
Thus \(\pi\) is a well-defined surjective algebra homomorphism.  Moreover,
\[
\pi((\Khat_i^-)^{-1}\Khat_i^+-1)=\Khat_i^{-1}\Khat_i-1=0,
\]
so \(J\subseteq\ker\pi\), and \(\pi\) factors through \(D/J\).

It remains to show that the induced map
\[
\overline\pi:D/J\longrightarrow U
\]
is injective.  In \(D/J\), the relations
\[
(\Khat_i^-)^{-1}\Khat_i^+=1
\]
imply
\[
\Khat_i^-=\Khat_i^+.
\]
Thus the two Cartan tori are identified.  Since the diagonal Cartan elements commute and are invertible, the ideal generated by \((\Khat_i^-)^{-1}\Khat_i^+-1\) is the same as the ideal generated by \(\Khat_i^--\Khat_i^+\).  This is precisely the standard quotient of the double of the two Borels which gives the full quantized enveloping algebra; the \(sl_2\) case is proved in
\href{https://staff.fnwi.uva.nl/j.v.stokman/quantumdouble2.pdf}{staff.fnwi.uva.nl}, Corollary~2.12, and the general case is the identical Drinfeld--Tanisaki argument.

One can also see injectivity directly from PBW bases.  The quotient \(D/J\) has a spanning set consisting of monomials
\[
F_{\beta_N}^{m_N}\cdots F_{\beta_1}^{m_1}\,
\Khat_1^{\lambda_1}\cdots \Khat_n^{\lambda_n}\,
E_{\beta_1}^{r_1}\cdots E_{\beta_N}^{r_N},
\]
with the single Cartan torus \(\Khat_1\cdots\Khat_n=1\).  The relations are exactly those of \(U_q(\sln)\), and the PBW theorem for \(U_q(\sln)\) says that the corresponding monomials form a basis of \(U\).  Since \(\overline\pi\) sends this spanning set to the PBW basis of \(U\), it is injective.

Therefore
\[
D/J\cong U_q(\sln).
\]
Finally, using Theorem~\ref{thm:main-isom},
\[
\bcS_n(\fB^*)/I
\cong
D/J
\cong
U_q(\sln).
\]
This completes the proof.
\end{proof}

\begin{remark}
Because \(\prod_i\Khat_i^+=1\) and \(\prod_i\Khat_i^-=1\), one of the \(n\) relations \((\Khat_i^-)^{-1}\Khat_i^+=1\) is redundant.  This is consistent with the fact that \(U_q(\sln)\) has rank \(n-1\).
\end{remark}

\begin{remark}
The choice \(\Phi(b_{ii})=(\Khat_i^-)^{-1}\) is forced by the right-directed convention for the horizontal arcs and by the form of the ideal \(b_{ii}c_{ii}-1\).  If one uses the opposite orientation convention for the \(b\)-arcs, the same statement appears with \(b_{ii}\) replaced by \(b_{ii}^{-1}\); the quotient relation is then \(b_{ii}^{-1}c_{ii}=1\).  With the convention in the question, the diagonal quotient is exactly \(b_{ii}c_{ii}=1\).
\end{remark}

\begin{thebibliography}{99}

\bibitem{KS}
A.~Klimyk and K.~Schm\"udgen,
\emph{Quantum Groups and Their Representations},
Springer, 1997.

\bibitem{Lusztig}
G.~Lusztig,
\emph{Introduction to Quantum Groups},
Progress in Mathematics, Birkh\"auser, 1993.

\bibitem{StokmanDouble}
J.~V.~Stokman,
\emph{Quantum double as a braided Hopf algebra and quantum \(sl_2\)},
lecture notes,
\href{https://staff.fnwi.uva.nl/j.v.stokman/quantumdouble2.pdf}{staff.fnwi.uva.nl}.

\bibitem{StokmanQuantumsl2}
J.~V.~Stokman,
\emph{Quantum \(sl(2)\) and the universal \(R\)-matrix},
lecture notes,
\href{https://staff.fnwi.uva.nl/j.v.stokman/Quantumsl2.pdf}{staff.fnwi.uva.nl}.

\bibitem{BorelPresentation}
\emph{Presentation by Borel subalgebras and Chevalley generators for quantum enveloping algebras},
\href{https://www.cambridge.org/core/services/aop-cambridge-core/content/view/886B0B4D8E3D242A51A0380A486351EA/S0013091504000689a.pdf/presentation-by-borel-subalgebras-and-chevalley-generators-for-quantum-enveloping-algebras.pdf}{cambridge.org}.

\bibitem{HalDouble}
\emph{The Drinfeld double: algebraic level},
\href{https://hal.science/hal-02159493v1/file/CQG.pdf}{hal.science}.

\bibitem{ArxivDoubleBorel}
\emph{Quantum double of the Borel subalgebra of \(U_q(sl_2)\)},
\href{https://arxiv.org/pdf/math/0512563}{arxiv.org}.

\bibitem{Turaev}
V.~G.~Turaev,
\emph{Quantum Invariants of Knots and 3-Manifolds},
de Gruyter, 1994.

\bibitem{Majid}
S.~Majid,
\emph{Foundations of Quantum Group Theory},
Cambridge University Press, 1995.

\end{thebibliography}

\end{document}